\chapter{A Probabilistic View of Learning}

\section{Randomness and Data Generation}

In Chapter 1, we introduced learning as the problem of approximating an unknown function from data. We now refine this perspective by incorporating \emph{uncertainty}.

\medskip

Data is not deterministic. Even when a true relationship exists, observations are subject to noise, variability, and incomplete information. A natural way to model this is through probability.

\medskip

\noindent
We assume that data is generated according to an unknown probability distribution
\[
P \quad \text{on} \quad \mathcal{X} \times \mathcal{Y}.
\]

A dataset
\[
\mathcal{D} = \{(x_i, y_i)\}_{i=1}^n
\]
is modeled as independent samples drawn from $P$:
\[
(x_i, y_i) \sim P.
\]

\medskip

\noindent
\textbf{Interpretation.}
\begin{itemize}
    \item The distribution $P$ captures all properties of the data-generating process.
    \item The dataset is only a finite, noisy observation of this process.
\end{itemize}

\medskip

This viewpoint shifts the goal of learning: instead of fitting a dataset, we aim to learn properties of the underlying distribution.

\section{Conditional Structure and Prediction}

A useful decomposition of $P$ is given by
\[
P(x,y) = P(y \mid x) P(x).
\]

\medskip

\noindent
\textbf{Key object:} the conditional distribution $P(y \mid x)$.

\medskip

\noindent
Prediction can be understood as estimating this conditional relationship.

\medskip

\noindent
\textbf{Example (Regression).}
\[
y = f^\star(x) + \varepsilon, \quad \mathbb{E}[\varepsilon \mid x] = 0.
\]
Then
\[
f^\star(x) = \mathbb{E}[y \mid x].
\]

\medskip

\noindent
\textbf{Example (Classification).}
\[
P(y = k \mid x), \quad k = 1, \dots, K.
\]
Prediction corresponds to choosing the most probable class:
\[
\hat{y} = \arg\max_k P(y = k \mid x).
\]

\medskip

Thus, many learning problems reduce to estimating conditional expectations or probabilities.

\section{Bayes' Theorem and Inference}

Bayes' theorem provides a fundamental relationship between conditional probabilities:
\[
P(y \mid x) = \frac{P(x \mid y) P(y)}{P(x)}.
\]

\medskip

\noindent
\textbf{Interpretation.}
\begin{itemize}
    \item $P(y)$ is the prior belief about $y$,
    \item $P(x \mid y)$ is the likelihood of observing $x$ given $y$,
    \item $P(y \mid x)$ is the posterior belief after observing $x$.
\end{itemize}

\medskip

In practice, Bayes' theorem suggests two complementary approaches:
\begin{itemize}
    \item \textbf{Discriminative modeling:} directly model $P(y \mid x)$
    \item \textbf{Generative modeling:} model $P(x \mid y)$ and $P(y)$
\end{itemize}

\medskip

\noindent
\textbf{Engineering perspective.}
\begin{itemize}
    \item Discriminative models often yield better predictive performance.
    \item Generative models provide richer structure and can generate data.
\end{itemize}

\section{Loss Functions as Statistical Objectives}

In Chapter 1, we introduced loss functions as measures of error. From a probabilistic viewpoint, loss functions arise naturally from statistical principles.

\medskip

\noindent
Consider a model $f_\theta$ that defines a conditional distribution $p_\theta(y \mid x)$.

\medskip

\noindent
\textbf{Negative log-likelihood.}
\[
\ell(f_\theta(x), y) = -\log p_\theta(y \mid x).
\]

\medskip

Minimizing the empirical risk
\[
\frac{1}{n} \sum_{i=1}^n -\log p_\theta(y_i \mid x_i)
\]
is equivalent to maximizing the likelihood:
\[
\max_\theta \prod_{i=1}^n p_\theta(y_i \mid x_i).
\]

\medskip

\noindent
\textbf{Examples.}
\begin{itemize}
    \item Gaussian noise $\Rightarrow$ squared loss
    \item Categorical distribution $\Rightarrow$ cross-entropy loss
\end{itemize}

\medskip

Thus, many commonly used loss functions are not arbitrary—they correspond to probabilistic assumptions about the data.

\section{Risk and Expected Loss}

The true objective of learning is to minimize the \emph{expected loss} (or \emph{risk})
\[
R(f) = \mathbb{E}_{(x,y) \sim P} \big[ \ell(f(x), y) \big].
\]

\medskip

\noindent
\textbf{Ideal problem.}
\[
\min_{f \in \mathcal{H}} R(f).
\]

\medskip

However, $P$ is unknown, so $R(f)$ cannot be computed exactly.

\medskip

\noindent
Instead, we approximate it using the empirical risk:
\[
\hat{R}_n(f) = \frac{1}{n} \sum_{i=1}^n \ell(f(x_i), y_i).
\]

\medskip

\noindent
\textbf{Empirical Risk Minimization (ERM).}
\[
\min_{f \in \mathcal{H}} \hat{R}_n(f).
\]

\medskip

ERM is the central principle underlying most machine learning algorithms.

\section{Approximation, Estimation, and Optimization}

The probabilistic formulation reveals three distinct sources of error:

\medskip

\noindent
\textbf{1. Approximation error}  
The model class $\mathcal{H}$ may not contain the true function.

\medskip

\noindent
\textbf{2. Estimation error}  
We only observe a finite sample from $P$.

\medskip

\noindent
\textbf{3. Optimization error}  
We may not find the exact minimizer of the empirical risk.

\medskip

\noindent
\textbf{Insight.}  
Modern machine learning systems are effective because they balance these three errors.

\section{Conditional Expectation as an Optimal Predictor}

A key theoretical result provides insight into the structure of optimal predictors.

\medskip

\noindent
\textbf{Proposition.}  
For squared loss, the function that minimizes the risk is
\[
f^\star(x) = \mathbb{E}[y \mid x].
\]

\medskip

\noindent
\textbf{Interpretation.}
\begin{itemize}
    \item The optimal prediction is the conditional mean.
    \item Learning aims to approximate this conditional expectation.
\end{itemize}

\medskip

\noindent
This result explains why regression models focus on estimating averages, even when data is noisy.

\section{From Probabilities to Algorithms}

While the probabilistic framework provides clarity, practical algorithms must operate under constraints:
\begin{itemize}
    \item limited data,
    \item high-dimensional inputs,
    \item computational limitations.
\end{itemize}

\medskip

\noindent
Thus, models are designed not only to be statistically sound, but also:
\begin{itemize}
    \item computationally efficient,
    \item numerically stable,
    \item scalable to large datasets.
\end{itemize}

\medskip

\noindent
This interplay between probability and computation is a defining feature of modern machine learning.

\section{A Unifying Perspective}

We can now refine the learning pipeline:

\begin{enumerate}
    \item Assume data is generated from an unknown distribution $P$
    \item Choose a model $p_\theta(y \mid x)$ or function $f_\theta$
    \item Define a loss derived from probabilistic principles
    \item Minimize empirical risk
\end{enumerate}

\medskip

\noindent
\textbf{Key takeaway.}  
Learning is the process of approximating conditional distributions from finite samples.

\section{Outlook}

This chapter introduced the probabilistic foundation of learning:
\begin{itemize}
    \item data as samples from a distribution,
    \item prediction as estimating conditional structure,
    \item loss functions as statistical objectives.
\end{itemize}

\medskip

In the next chapter, we study \emph{hypothesis spaces} and \emph{inductive bias}, which determine what kinds of functions can be learned and how generalization becomes possible.

\medskip

\noindent
\textbf{Guiding principle.}  
Probability provides the language of uncertainty, but successful learning requires combining it with structure, computation, and approximation.
